\section{Reading the profile in concrete families}\label{sec:examples}

The exponential term alone does not distinguish all the regimes of the
answer.  The following families separate three further effects: the
comparison of named suffixes, the count associated with leaving an owner,
and the ordinary terminal constraint.  They also show how several
comparisons combine, first within one owner and then across different
owners.  Every partition into owners below is the partition of the actual
maximizer of \eqref{profile:saddle}.

Throughout this section, \(\log3+L\) means coordinatewise addition.  A
length vector specifies actual arithmetic cutoffs by
\[
 y_0=3,\qquad \log y_i=e^{L_i}\log y_{i-1}.
\]
Consequently every displayed estimate for \(\Phi_k(\log3+L)\) is also
an estimate for \(H^{(k+1)}(x,\mathbf y,2\mathbf y)/x\) throughout the
admissible region in Theorem~\ref{paper:main-H}.  Taking \(N_i=2y_i\)
for \(i\le k\) and \(N_{k+1}\ge N_k\) gives the corresponding
multiplication-table estimate by Theorem~\ref{paper:main-table}.
The dimension is fixed in all comparisons.  A constant denoted by \(C\)
is chosen once, sufficiently large in terms of that dimension.

For convenient reference, write the correlation in
\eqref{profile:edge-factor} as
\begin{equation}
 \rho(r,q;z)=
 \frac{\{\log(r+z)-p\log(q+z)\}^2}
 {\{\log(r+z)-p\log(q+z)\}^2+
             \{\log(q+z)\}^2p(1-p)},\qquad
 p=\frac{q+z}{r+z}.
 \label{ex:rho}
\end{equation}
Here \(r,q,z\) retain their meanings in the profile; in particular,
\(p\) is not the correlation.  The proofs of the evaluations, including
the finite inequalities used to locate the constants, are given in
Appendix~\ref{app:examples}.

\subsection{Equal cutoffs}

Put \(H=k+1\) and
\begin{equation}
 \alpha_H=\frac{\log((H-1)/\log H)}{\log H},\qquad
 I_H=1+(H-1)\alpha_H-H^{\alpha_H}.
 \label{ex:equal-rate}
\end{equation}
For \(L=(T,0,\ldots,0)\), the cutoffs are all \(3^{e^T}\), and
\begin{equation}
 \Phi_k(\log3+L)\asymp_k e^{-I_HT}T^{-3/2}.
 \label{ex:equal-answer}
\end{equation}
To see the factors in this formula, replace the zero lengths by a fixed
\(C\).  The resulting saddle has one terminal owner, with common product
\(\alpha_H+O_k(T^{-1})\).  All binary factors have a positive constant
lower bound.  The ordinary suffix starts in the first group, has zero
total drift, and has one group of length \(T\), followed by bounded
total length.  Its root factor is \(T^{-1/2}\) and its ballot factor is
\(T^{-1}\).  Bounded arithmetic stability transfers the estimate back
to the original equal cutoffs.  In particular the zero coordinates have
not removed any divisor coordinate.

\subsection{Three six-dimensional regimes}

The next two families have a long first group, four bounded groups, and
a long last group.  Their different relative lengths expose different
parts of the answer.

First define
\begin{gather}
 \sigma_c=1-\frac{\log((7\log7-2\log2)/5)}{\log7},\qquad
 e_h=h^{\sigma_c}\log h-(h-1),
 \label{ex:critical-constants}\\
 a=-\frac{e_7}{e_2},\qquad
 b=-\frac{e_6+e_5+e_4+e_3}{e_2},\qquad
 L=(T,C,C,C,C,aT+bC).
 \label{ex:critical-lengths}
\end{gather}
The constants satisfy \(\alpha_2<\sigma_c<\alpha_3\), \(54<a<55\),
and \(66<b<67\).  The unique saddle is exactly
\(s_1=\cdots=s_6=\sigma_c\).  Nevertheless a binary comparison remains
unsaturated.  It is the first-group edge \(6\to1\), which keeps name
\(6\) and removes names \(1,\ldots,5\).  Put
\[
 \Delta=\sum_{h=3}^6
 \left\{h^{\sigma_c-1}(h\log h-2\log2)-(h-2)\right\}>0.
\]
The pair giving the smallest endpoint scales has cuts
\[
\begin{array}{c|cc}
 &\tau_1,\ldots,\tau_5&\tau_6\\ \hline
 \text{left}&0,\ldots,0&V_6\\
 \text{right}&V_1,\ldots,V_1&V_6
\end{array}
\]
and both reference costs equal \(C\Delta\).  Both completions are late.
On the left the first group is good, and on the right it is not; the four
middle groups are good on both sides.  Thus the full scores are
\begin{equation}
 X=1+C\Delta+\sqrt T+4\sqrt C,\qquad
 Y=1+C\Delta+4\sqrt C
 \quad (T\text{ sufficiently large}).
 \label{ex:critical-endpoints}
\end{equation}
The distinction between reference cost and named endpoint score changes
the power of \(T\).  With
\(\chi_c=\chi(\sigma_c,\rho(6,1;1))\), the answer is
\begin{equation}
 \Phi_6(\log3+L)\asymp_6
       e^{-J(L)}T^{-1/2-\chi_c}.
 \label{ex:critical-answer}
\end{equation}
There is no proper-owner count.  The ordinary suffix consists of the
last group and has a drift proportional to its length, leaving only
the power \(T^{-1/2}\).  All remaining binary factors are bounded below.
The exponent is explicit: if
\(I_h^{c}=1+(h-1)\sigma_c-h^{\sigma_c}\), then
\[
 J(L)=(I_7^{c}+aI_2^{c})T+
             (I_6^{c}+I_5^{c}+I_4^{c}+I_3^{c}+bI_2^{c})C.
\]

Now move the last cutoff much farther away by taking
\begin{equation}
 L=(T,0,0,0,0,T^3).
 \label{ex:split-lengths}
\end{equation}
Write \(\alpha=\alpha_2\).  There is a unique \(s_*\in(\alpha,1)\)
satisfying
\begin{equation}
 z_*=2^{\alpha/s_*},\qquad
 (5+z_*)^{s_*-1}
 \{(5+z_*)\log(5+z_*)-z_*\log z_*\}=5;
 \qquad 0.54070<s_*<0.54072.
 \label{ex:split-root}
\end{equation}
After replacing the four zero lengths by \(C\), the actual owners are
\([1,5]\) and \([6]\).  Their products tend to \(s_*\) and \(\alpha\),
respectively.  Only the proper edge \(5\to0\) in the first group is
unsaturated.  Its endpoint scales are both bounded, giving
\(T^{-2\chi_*}\), where
\(\chi_*=\chi(s_*,\rho(5,0;z_*))\).  The proper count gives
\(T^{-s_*/2}\).  The terminal length is \(T^3\), but its drift has order
\(T\), supplied by the earlier owner.  Hence its joint ordinary and root
factor is \(T(T^3)^{-3/2}\).  Altogether, with
\[
 I_2=1+\alpha-2^\alpha,\qquad
 I_{\rm p}=1+5s_*+\alpha-(5+z_*)^{s_*},
\]
we obtain
\begin{equation}
 \Phi_6(\log3+L)\asymp_6
 e^{-I_2T^3-I_{\rm p}T}
 T^{1-s_*/2-2\chi_*}(T^3)^{-3/2}.
 \label{ex:split-answer}
\end{equation}
The child in this proper comparison has \(z_*>1\), so \(q=0\) still
leaves a nondegenerate binary mark.  It is distinct from the terminal
ordinary comparison.

There is a precise dimensional reason for using six coordinates in
these two families.  For the first construction, replacing \(7\) by
\(H\) in \eqref{ex:critical-constants} gives a product below
\(\alpha_2\) for \(3\le H\le6\), and above it for \(H=7\).  For the
second, the boundary sign is
\[
 f(H)=H^{\alpha-1}(H\log H-2\log2)-(H-2),\qquad
 f(6)>0>f(7).
\]
Strict concavity of \(f\) on \([2,\infty)\) makes the strict branch
possible precisely for \(H\ge7\) in this two-scale family.  These are
thresholds for the stated constructions, not classifications of all
parameter vectors.

A final contrast isolates the proper count.  Fix
\(s=107/200\), \(\theta=267/500\), and \(z=2^{\theta/s}\).  Set
\begin{gather}
 f_r=(r+z)^{s-1}E_r(z)-r,\qquad
 c_r=1-z\log2\,(r+z)^{s-1},\nonumber\\
 a=-\frac{f_1+f_2+f_3+f_4}{f_5},\qquad
 b=\frac{ac_5+c_1+c_2+c_3+c_4}{2^\theta\log2-1},\qquad
 L=t(a,1,1,1,1,b).
 \label{ex:ray-lengths}
\end{gather}
Here \(a\) and \(b\) are positive (approximately \(10.958044\) and
\(1651.936051\)).  The exact owners remain \([1,5]\), \([6]\), with
products \(s,\theta\), but every binary factor is bounded below.
Writing \(J_0=J(a,1,1,1,1,b)\), we have
\begin{equation}
 \Phi_6(\log3+L)\asymp_6 e^{-tJ_0}t^{-(1+s)/2}.
 \label{ex:ray-answer}
\end{equation}
Thus leaving an owner can impose a count penalty even when every binary
comparison has sufficient endpoint room.

\subsection{Two comparisons in one owner}

The next family has two unsaturated edges in the same proper owner.
They share its count factor, but have different endpoint scales and
different correlations.

For \(x>y\ge1\), define
\[
 t(x,y)=1-\frac{\log((x\log x-y\log y)/(x-y))}{\log x}.
\]
Let \(z\) be the unique solution in \((1.98455,1.98457)\) of
\begin{equation}
 t(14+z,6+z)=t(6+z,z),\qquad
 \sigma=t(6+z,z),\qquad \theta=\frac{\sigma\log z}{\log2}.
 \label{ex:same-root}
\end{equation}
Then \(\alpha_2<\theta<\sigma\), with
\(\sigma\) approximately \(0.551693\) and \(\theta\) approximately
\(0.545525\).  For \(1\le r\le14\), put
\[
 \kappa_r=(r+z)^{\sigma-1},\qquad
 f_r(q)=q-\kappa_r E_q(z),\qquad g_r=f_r(r).
\]
The defining equations give the exact identities
\begin{equation}
 g_6=0,\qquad f_{14}(6)=g_{14}>0,
 \qquad g_1<0.
 \label{ex:same-balances}
\end{equation}
Take \(k=15\), \(U=T^3\), and specify the lengths by
\begin{align}
 &L_1=T,\quad L_9=U,\quad
 L_i=C\quad(2\le i\le13,\ i\ne9),\nonumber\\
 &L_{14}=W=-\frac{g_{14}T+
                C\sum_{\substack{2\le r\le13\\r\ne6}}g_r}{g_1},
 \qquad
 L_{15}=R=
 \frac{\sum_{i=1}^{14}L_i(1-z\log2\,\kappa_{15-i})}
 {2^\theta\log2-1}.
 \label{ex:same-lengths}
\end{align}
All these lengths are positive; \(W\asymp T\) and \(R\asymp U\).
The actual saddle is
\((\sigma,\ldots,\sigma,\theta)\), so the owners are \([1,14]\) and
\([15]\).  The two equations defining \(W,R\) balance the common
proper-owner direction and the terminal direction.  The remaining KKT
inequalities are strict.

The first rare comparison is \(14\to6\) in group \(1\).  Keeping the
six deepest names until their deadlines gives equal bounded reference
costs at its endpoints.  The left completion makes the first group good,
while the right completion does not.  The second rare comparison is
\(6\to0\) in group \(9\).  Its six names must survive through the
earlier first group, and this creates a cost of order \(T\) at both
endpoints.  Precisely,
\begin{equation}
 (X_{1,6},Y_{1,6})\asymp(\sqrt T,1),\qquad
 (X_{9,0},Y_{9,0})\asymp(T,T).
 \label{ex:same-endpoints}
\end{equation}
Every other binary factor is bounded below.  Put
\[
 \chi_1=\chi(\sigma,\rho(14,6;z)),\qquad
 \chi_2=\chi(\sigma,\rho(6,0;z)).
\]
There is one proper count, of order \(U^{-\sigma/2}\), and one ordinary
root factor, of order \(U^{-1/2}\).  The resulting full profile is
\begin{equation}
 \Phi_{15}(\log3+L)\asymp_{15}
 e^{-J(L)}U^{-(1+\sigma)/2}T^{-\chi_1}
                   \left(\frac{T^2}{U}\right)^{2\chi_2},
 \qquad U=T^3.
 \label{ex:same-answer}
\end{equation}
Its exponential is explicit as well:
\begin{equation}
 J(L)=\sum_{i=1}^{14}L_i
  \{1+(15-i)\sigma+\theta-(15-i+z)^\sigma\}
       +R(1+\theta-2^\theta).
 \label{ex:same-rate}
\end{equation}
This example explains why the profile takes the product of all canonical
edge factors inside an owner, but assigns a count factor only once to
that owner.  No assertion of independence between the two comparisons
is involved.

\subsection{Two comparisons in different owners}

One more family makes the different owner powers visible.  Put
\(k=18\), \(\alpha=\alpha_2\), and, for \(c>0\), define
\[
 F(s;c)=(1+e^{c/s})^{s-1}
       \{(1+e^{c/s})\log(1+e^{c/s})-e^{c/s}(c/s)\}-1.
\]
Let \(v,u\) be the roots
\begin{equation}
 F(v;\alpha\log17)=0,\qquad
 G_2^*=v\log(1+17^{\alpha/v}),\qquad F(u;G_2^*)=0.
 \label{ex:multi-roots}
\end{equation}
They are uniquely determined in \((0,1)\) and satisfy
\(u>v>\alpha\); approximately they are
\(u=0.534353736\), \(v=0.531542356\).
Set
\begin{gather}
 z_2^*=17^{\alpha/v},\quad z_1^*=e^{G_2^*/u},\quad
 Z_j^*=(1+z_j^*)^{s_j^*},\quad (s_1^*,s_2^*)=(u,v),\nonumber\\
 L_1=T,\quad L_2=T^3,\quad
 L_i=2^{i-3}T^5\ (3\le i\le17),\quad L_{18}=T^{11}.
 \label{ex:multi-lengths}
\end{gather}
For sufficiently large \(T\), the exact owners are
\([1]\), \([2]\), \([3,18]\).  Their products satisfy
\[
 u_T=u+O(T^{-2}),\qquad v_T=v+O(T^{-2}),\qquad
 \tau_T=\alpha+O(T^{-6}).
\]
Each of the first two owners has just the proper edge \(1\to0\).
Their endpoint scales are
\[
 (X_1,Y_1)\asymp(1,1),\qquad
 (X_2,Y_2)\asymp(T,1).
\]
The old first group supplies the left endpoint room in the second
comparison.  With
\(\chi_1=\chi(u,\rho(1,0;z_1^*))\) and
\(\chi_2=\chi(v,\rho(1,0;z_2^*))\), the respective binary costs are
\(T^{-2\chi_1}\), \(T^{-4\chi_2}\), and the respective count costs
are \(T^{-u/2}\), \(T^{-3v/2}\).

All \(120\) terminal binary factors have a positive constant lower
bound.  The ordinary suffix is the last group: its length is \(T^{11}\)
and its drift has order \(T^5\).  Its single joint factor is therefore
\(T^{-23/2}\).  Define
\begin{align}
 J_0(T)={}&T(1+u+v+16\alpha-Z_1^*)
       +T^3(1+v+16\alpha-Z_2^*)\nonumber\\
 &+T^5\sum_{i=3}^{17}2^{i-3}
       \{1+(h_i-1)\alpha-h_i^\alpha\}
       +T^{11}(1+\alpha-2^\alpha),\qquad h_i=20-i.
 \label{ex:multi-rate}
\end{align}
Then \(J(L)=J_0(T)+O(T^{-1})\), and
\begin{equation}
 \Phi_{18}(\log3+L)\asymp_{18}
 e^{-J_0(T)}
 T^{-23/2-u/2-3v/2-2\chi_1-4\chi_2}.
 \label{ex:multi-answer}
\end{equation}
Thus the two owners contribute their own powers and counts, while the
whole chain still has only one ordinary root.

For this particular construction the dimensional threshold is governed
by a one-step, rather than a long-interval, secant.  If
\[
 t_1(x)=1-\frac{\log\{x\log x-(x-1)\log(x-1)\}}{\log x},
\]
then \(t_1(17)<\alpha_2<t_1(18)\).  Together with the one-valley
property in Lemma~\ref{hier:lem:valley}, this says that a terminal child
with limiting product \(\alpha_2\) needs at least \(17\) names before
the adjacent singleton can have a strictly larger product.  Two such
singleton owners therefore first fit in dimension \(18\).  This
statement concerns the specified singleton construction; the
fifteen-dimensional example uses a different arrangement of ranks.
